\paragraph{R4-O5, R4-W3 (Baselines are general retrieval/indexing; experiments show retrieval efficiency, not counterfactual capability).}

\begin{itemize}[leftmargin=1.2em, itemsep=1pt, topsep=2pt]
\item \textit{O5: ``The current baselines (e.g., DuckDB, KD-Tree, BM25) are mainly general retrieval/indexing methods. Therefore, the experiments mainly demonstrate retrieval efficiency rather than counterfactual capability. The authors are encouraged to include more targeted counterfactual baselines, or discuss why such comparisons are difficult.''}
\item \textit{W3 (part): ``\dots evaluation baselines would benefit from more detailed justification and discussion.''}
\end{itemize}
\smallskip

\noindent\textbf{Response.}
We agree that the original baselines measure retrieval efficiency, not counterfactual capability, and we add a capability-oriented evaluation. Table~\ref{tab:baselines} reports a controlled simulation comparing DcCA against reference reimplementations of representative counterfactual generators---a data-grounded feasible-path method (FACE~\cite{poyiadzi2020face}), a diverse gradient method (DiCE~\cite{mothilal2020explaining}), an unconstrained gradient method (Wachter~\cite{wachter2017counterfactual}), and a genetic search (GeCo~\cite{schleich2021geco})---on validity, feasibility (immutability), in-distribution rate, set diversity, cost, and runtime. The pattern is clear: data-grounded methods are in-distribution by construction, but only DcCA also respects immutability (via immutability-matched retrieval), while model-centric generators achieve low nominal cost yet are infeasible ($0\%$) and off-manifold. DcCA also yields the most diverse valid set. We retain the retrieval/indexing baselines (DuckDB, KD-Tree, BM25) to characterize the efficiency dimension.

\begin{table}[h]
\centering\small
\caption{Counterfactual capability: DcCA vs.\ model-centric generators. \emph{Controlled simulation on synthetic tabular data ($N{=}6000$, $3$ categorical immutable attributes, $k{=}5$, $150$ queries) with reference reimplementations of each method's core mechanism---not the released libraries and not UCI data. Full-library results (incl.\ DACE~\cite{kanamori2020dace}, prototype-CF~\cite{vanlooveren2020interpretable}) on Adult, German-Credit, and HELOC are to follow. \bb{Are we doing additional experiments? If not we should delete this}} Higher is better except Cost.}
\label{tab:baselines}
\begin{tabular}{lrrrrr}
\toprule
Method & Valid.\% & Feas.\% & In-dist.\% & $\Omega_{\text{div}}$ & Cost \\
\midrule
\textbf{DcCA (data-centric)}        & 100.0 & \textbf{64.7} & \textbf{100.0} & \textbf{27.08} & 4.09 \\
FACE (data-path)~\cite{poyiadzi2020face}      & 100.0 & 1.9  & 100.0 & 17.29 & 2.83 \\
DiCE (diverse grad)~\cite{mothilal2020explaining}   & 100.0 & 0.0  & 1.1   & 9.32  & \textbf{1.73} \\
Wachter (grad)~\cite{wachter2017counterfactual}        & 100.0 & 0.0  & 0.0   & 2.27  & 1.89 \\
GeCo (genetic)~\cite{schleich2021geco}        & 80.8  & 0.0  & 0.7   & 20.24 & 1.77 \\
\bottomrule
\end{tabular}
\end{table}

\noindent\emph{Revision (Section~6 summary; details in Appendix~A):}
\begin{quote}
\textcolor{blue}{Beyond runtime, we further study counterfactual \emph{quality}---collective-validity objectives, a comparison against representative counterfactual generators (FACE, DiCE, Wachter, GeCo), and approximation/scalability---together with the practical index-construction and maintenance costs; for space, these are reported in Appendix~A. In brief, objective-specific set selection improves every collective-validity objective over top-$k$ (which is best for none), and \textbf{DcCA} uniquely returns feasible, in-distribution counterfactuals whereas model-centric generators are infeasible and off-manifold.}
\end{quote}
